\section*{致谢}

我们感谢 S. Petrarca 在本研究早期阶段的参与以及许多讨论。我们还感谢 E. Gabrielli 与 LPTHE
的成员们的讨论。G.M. 与 C.P. 感谢 CERN 理论部，本研究的部分工作是在那里完成的。我们感谢
M.U.R.S.T. 与欧共体合同 CHRX-CT92-0051 的部分资助。C.P. 感谢 Human Capital and Mobility
计划（合同 ERBCHBICT930887）的支持。CTS 感谢 Particle Physics and Astronomy Research Council
通过 Senior Fellowship 的支持。

\begin{figure}[htbp]
\centering
\includegraphics[width=0.62\textwidth]{images/file.pdf}
\caption{标准夸克（a）与标量夸克（b）的图。}
\label{fig:appe}
\end{figure}

\section*{附录}

在本附录中我们证明：对前向两夸克矩阵元的非微扰贡献被压低为外动量 $p$ 的 $p^2$ 逆幂。考虑非截腿
两费米子格林函数
\beq F_O(p)=\int d^4x d^4y e^{-i p \cdot x} \langle \bar q(0) \Gamma
q(x) O_\Gamma(y) \rangle ,\nn \eeq
其中 $O_\Gamma(y)=\bar q(y) \Gamma q(y)$，$\Gamma$
是某个狄拉克矩阵。自然会担心 $F_O(p)$ 的红外行为，因为它包含零动量算符的插入。一般地，
我们预期红外发散会出现在以例外动量计算的格林函数中。用物理语言说，这是因为零动量的
$O_\Gamma$ 矩阵元携带着关于物理质量谱的信息，而这是微扰论无法触及的。特别地，若算符
$O_\Gamma$ 是赝标密度（$\gamma_5$）且我们在手征极限下工作，则对应于 π 介子的 Goldstone 极点贡献应当
出现在 $F_O(p)$ 中。$F_O(p)$ 中非微扰贡献存在的标志是微扰展开中出现红外发散。一般情形正是如此。
然而对于大 $p^2$ 极限下的夸克自由度，算符乘积展开（OPE）保证 $F_O(p)$ 可以在微扰论所有阶可靠地
计算。论证如下。在大 $p^2$ 极限下，$F_O(p)$ 的主导贡献来自被积函数奇异的积分区域（关于 $x$ 与
$y$）。这意味着在此极限下我们必须研究
\beq F_O(x,y)=\langle \bar q(0) \Gamma q(x) O_\Gamma(y)  \rangle,\nn \eeq
在 $x \sim 0$ 与 $x \sim y$ 处的行为。当然在这两种情形都可以使用如下形式的
OPE\footnote{从技术角度看，OPE 是一个弱算符陈述；当 $O_\Gamma$ 插入含其他算符的格林函数时，
应使用略微修改的形式。不过容易确信，一般情形所需的附加项与本讨论无关。}：
对 $x \sim y$，
\beq q_\alpha (x) O_\Gamma (y) \to A^\Gamma_{\alpha\beta}(x-y)
 \Gamma_{\beta\delta} q_\delta (y)+\dots, \nn \eeq
对 $x \sim 0$，
\beq \bar q(x) \Gamma q(0) \to B_\Gamma (x^2) \bar q(0) \Gamma q(0) +\dots
,\nn \eeq

由于 QCD 渐近自由，$A^\Gamma_{\alpha\beta}(x-y) $ 与 $B_\Gamma (x)$ 的行为可以在微扰论中计算。
相互作用对自由场行为给出对数修正：
\beq  A_\Gamma (x) =\xslash C_\Gamma (x^2) \equiv \xslash \frac{a_\Gamma
\ln^{\alpha_\Gamma}(\mu^2 x^2)+\dots}{x^2} \nn \eeq
\beq B_\Gamma(x^2)= b_\Gamma \ln^{\beta_\Gamma}(\mu^2 x^2)+\dots \nn \eeq
利用上式，在 $p^2 \to \infty$ 极限下我们发现
\beq \int d^4x d^4y e^{-i p \cdot x} \left( \langle \bar q(0)\Gamma  A_\Gamma
 (x-y) \Gamma
q(y) \rangle + B_\Gamma(x^2) \langle \bar q(0) \Gamma q(0) O_\Gamma(y) \rangle
 \right) = \nn \eeq \beq -i p_\mu \bar C_\Gamma (p^2) \bar S^\mu_\Gamma (p)+
\bar B_\Gamma (p^2) \bar \Delta_\Gamma (0) ,\nn \eeq
其中
\beq \bar S^\mu_\Gamma (p)= \int d^4y e^{-i p \cdot y}\langle \bar q(0)
 \Gamma
\gamma^\mu \Gamma q(y) \rangle
\,\,\,\,\,\,\,\,\,\,\,\,\,
 \bar \Delta_\Gamma (0)= \int d^4y \langle \bar q(0) \Gamma q(0)
O_\Gamma(y) \rangle . \nn \eeq

项 $-i p_\mu \bar C_\Gamma (p^2) \bar S^\mu_\Gamma (p)$ 可在微扰论中计算。按量纲分析，当 $p^2 \to \infty$
时其行为为
\beq -i p_\mu \bar C_\Gamma (p^2) \bar S^\mu_\Gamma (p) \to
c_\Gamma \frac{\ln^{\gamma_\Gamma}(p^2/\mu^2)}{p^2} . \nn \eeq
至于 $\bar B_\Gamma (p^2)$，我们求得
\beq   \bar B_\Gamma (p^2) \to
d_\Gamma \frac{\ln^{\delta_\Gamma}(p^2/\mu^2)}{p^4} . \nn \eeq
但 $\bar B_\Gamma (p^2)$ 被非微扰量 $\bar \Delta_\Gamma (0)$ 相乘。然而我们已经证明：在大 $p^2$
极限下，微扰贡献比非微扰贡献多压低一个 $p^2$ 幂次。这就是为什么即使在外动量为例外动量时，
只要 $p^2$ 很大，$F_O(p)$ 在微扰论中也是红外安全的。通过分析相应的费曼图可以直截了当地追踪这一
红外发散的压低。例如图 \ref{fig:appe}-a 中图的红外发散部分因下方胶子线中的动量流而被压低一个
$p^2$ 幂次。

考察在（假想的）以标量场表示夸克的理论中会发生什么是有趣的。此时红外发散不被压低，如图
\ref{fig:appe}-b 所示。一般论证如下：
\beq  F(p^2) =\int d^4x d^4y e^{-i p \cdot x} \langle \bar \phi (0) \phi(x)
J(y) \rangle, \nn \eeq  其中
\beq J(y)= \bar \phi (y) \phi(y). \nn \eeq
当 $x \sim 0$ 与 $x \sim y$ 时，OPE 给出
\beq  \phi (x) J(y) \to A\left((x-y)^2\right) \phi (y) + \dots
\,\,\,\,\,\,\,\,\,\,\,\,\,
 \bar \phi (x) \phi (0) \to B(x^2) \bar \phi (0) \phi (0) + \dots, \nn \eeq
其中
\beq  A (x^2) =  \frac{a
\ln^{\alpha}(\mu^2 x^2)+\dots}{x^2}
\,\,\,\,\,\,\,\,\,\,\,\,\,
B(x^2)= b \ln^{\beta}(\mu^2 x^2)+\dots, \nn \eeq
于是将得到
\beq F(p^2) \to \bar A(p^2) \bar S(p^2) + \bar B(p^2) \bar \Delta(0) \nn \eeq
其中
\beq \bar S (p^2)= \int d^4y e^{-i p \cdot y}\langle \bar \phi(0)
 \phi(y) \rangle
\,\,\,\,\,\,\,\,\,\,\,\,\,
 \bar \Delta (0)= \int d^4y \langle \bar \phi(0)  \phi(0)
J(y) \rangle . \nn \eeq
现在得到
\beq  \bar A(p^2)  \bar S (p^2) \to
c \frac{\ln^{\gamma}(p^2/\mu^2)}{p^4}  \nn \eeq
及
\beq   \bar B (p^2) \to
d \frac{\ln^{\delta}(p^2/\mu^2)}{p^4} . \nn \eeq
后两式意味着非微扰贡献与微扰贡献具有相同的 $p^2$ 依赖性。因此我们预期 $F(p^2)$ 中会出现红外发散。

\begin{thebibliography}{99}
\bibitem{bo} M. Bochicchio, L. Maiani, G. Martinelli, G.C. Rossi
and M. Testa, Nucl. Phys. B262 (1985) 331.
\bibitem{ma} L. Maiani and G. Martinelli, Phys. Lett. B178 (1986) 265.
\bibitem{wi} G. Martinelli, S. Petrarca, C.T. Sachrajda
and A. Vladikas, Phys. Lett. B311 (1993) 241.
Rather than quoting the numerical results of ref.[3], obtained
on a sample of 18 gauge configurations, we have opted for those
of the updated analysis of ref.\ \cite{gribov}, done with 36 gauge
configurations.
\bibitem{dallas} G. Martinelli et al.,
presented by C.T. Sachrajda, Nucl. Phys.B (Proc. Suppl.)
34 (1994) 507.
\bibitem{ber1} C. Bernard et al.,
Nucl. Phys. B (Proc. Suppl.) 17 (1990) 593; Nucl. Phys. B (Proc. Suppl.)
20 (1991) 410.
\bibitem{te} L. Maiani et al.,
Phys. Lett. B176 (1986) 445; Nucl. Phys. B289 (1987) 505
\bibitem{gav1} M.B. Gavela et al.,
Phys. Lett. B211 (1988) 139.
\bibitem{mms} L. Maiani, G. Martinelli and C.T. Sachrajda,
Nucl. Phys. B368 (1992) 281.
\bibitem{martsukuba} G. Martinelli,
 Nucl. Phys. B (Proc. Suppl.) 26 (1992) 31.
\bibitem{lepage}
G.P. Lepage and P.B. Mackenzie, Nucl. Phys. B (Proc. Suppl.) 20
(1991) 173; Phys. Rev. D48 (1993) 2250.
\bibitem{effi} M.A. Shifman and M.B. Voloshin,
Sov. J. Nucl. Phys. 45 (1987) 292; 47 (1988) 511.
\bibitem{powi} H.D. Politzer and M.B. Wise,
\pl{B206} (1988) 681; \pl{208} (1988) 504.
\bibitem{eh} E. Eichten and B. Hill, \pl {B240} (1990) 193.
\bibitem{bg} B. Grinstein, Nucl. Phys. B 339 (1990) 253.
\bibitem{efff} H. Georgi, Phys. Lett. B240 (1990) 447.
\bibitem{slami} B. Sheikholeslami and R. Wohlert, Nucl. Phys.
B259 (1985) 572.
\bibitem{altarelli} G. Altarelli, G. Curci, G.Martinelli
and S. Petrarca, Nucl. Phys. B187 (1981) 461.
\bibitem{buras} A.J. Buras, M. Jamin, M.E. Lautenbacher and P.H. Weisz,
Nucl. Phys.  B370 (1992) 69, Addendum ibid. Nucl. Phys.  B375
(1992) 501;
\noindent
A.J. Buras, M. Jamin and M.E. Lautenbacher, Nucl. Phys.  B400 (1993) 37 and
B40.
\bibitem{ciu} M. Ciuchini, E. Franco, G. Martinelli and L. Reina,
Phys. Lett. B301 (1993) 263; Nucl. Phys.  B415 (1994) 403.
\bibitem{altri} C. Parrinello, S. Petrarca and
A. Vladikas, Phys. Lett. B268 (1991) 236;
M. Paciello et al.,
Phys. Lett. B276 (1992) 163; Phys. Lett. B281 (1992) 417;
Phys. Lett. B289 (1992) 405.
\bibitem{gribov} M. Paciello et al., Rome prep. 1034 (July 1994),
to appear in Phys. Lett. B.
\bibitem{ka} L.H. Karsten and J. Smit, Nucl. Phys. B183
(1981) 103.
\bibitem{4ferm}   G. Martinelli, Phys. Lett. B141 (1984) 395; \\
C. Bernard, T. Draper and A. Soni,
 Phys. Rev. D36 (1987) 3224; \\
R. Frezzotti, E. Gabrielli, C. Pittori and
G.C. Rossi, Nucl. Phys. B373 (1991) 781;
Nucl. Phys. B409 (1993) 382.
\bibitem{msv} G. Martinelli, C.T. Sachrajda
and A. Vladikas, Nucl. Phys. B358 (1991) 212.
\bibitem{noi} G. Heatlie, G. Martinelli, C. Pittori,
G.C. Rossi and C.T. Sachrajda, Nucl. Phys. B352 (1991) 266.
\bibitem{tass2} G. Martinelli, C.T. Sachrajda, G. Salina
and A. Vladikas, Nucl. Phys. B378 (1992) 591.
\bibitem{newi} A. Borrelli, R. Frezzotti, E. Gabrielli and C. Pittori,
Nucl. Phys. B409 (1993) 382
\bibitem{hv} 't Hooft and M. Veltman, Nucl. Phys. B44 (1972) 189.
\bibitem{2ferm}
G. Martinelli and Y.C. Zhang, Phys. Lett. B123 (1983) 433; \\
E. Gabrielli, G. Heatlie, G. Martinelli, C. Pittori and
C.T. Sachrajda, Nucl. Phys. B362 (1991) 475.
\bibitem{giusti} L. Giusti and M. Testa, in preparation.
\bibitem{curci} G. Curci, Phys. Lett. B167 (1986) 425.

\end{thebibliography}
